\documentclass[10pt,a4paper]{amsart}
\usepackage[margin=19mm]{geometry}
\usepackage{amsmath,amssymb,mathtools}
\usepackage[scr=rsfs,cal=euler]{mathalfa}
\usepackage{xfrac}
\usepackage[unicode,hidelinks,bookmarks=false]{hyperref}
\newcommand{\ie}{i.e.\ }
\newcommand{\cf}{cf.\ }
\newcommand{\resp}{respectively\ }
\newcommand{\Subsection}{\subsection}
\providecommand{\nobreakdash}{\nobreak\hbox{-}}
\setlength{\emergencystretch}{2em}
\multlinegap0pt
\usepackage{bm}
\usepackage{bbm}
\usepackage{dsfont}
\usepackage{xspace}
\usepackage{cancel}
\usepackage{mathtools}
\usepackage[shortlabels]{enumitem}
\usepackage{amsthm}
\makeatletter
\@ifundefined{lemma}{\newtheorem{lemma}{Lemma}}{}
\@ifundefined{proposition}{\newtheorem{proposition}{Proposition}}{}
\makeatother
\title[Global in Time Estimates for Multi-phase Muskat Problem]{Global in Time Estimates for Multi-phase Muskat Problem}
\author{Zirui Wang}
\date{Source: arXiv:2604.12785v2 (CC BY 4.0)}
\begin{document}
\maketitle

\noindent\emph{Source and licence.} Global in Time Estimates for Multi-phase Muskat Problem. Version arXiv:2604.12785v2 (CC BY 4.0), distributed under the Creative Commons Attribution 4.0 International licence, as recorded in the arXiv licence metadata for this exact version and shown on the abstract page https://arxiv.org/abs/2604.12785v2. Full rights notice: licenses/arxiv-2604.12785-CC-BY-4.0.txt.\par\smallskip

\noindent\textit{Editorial scope.} Counted unit: the Proposition whose source label is \texttt{prop linear matrix} in the article (source0.tex lines 329--349, source label \texttt{prop linear matrix}), delivered together with the article's own complete original proof of it (source0.tex lines 350--425, one \texttt{proof} environment of 76 lines). No mathematics is written, repaired or re-derived: statement and proof are the article's own text line for line. Required context retained uncounted: Bochner lemma (lines 311--317). Every definition, display and label of those lines is retained unchanged. Omitted from this package and not redistributed: 1--310, 318--328, 426--812. Nothing inside a retained excerpt is omitted or altered: every retained body is delivered exactly as the article's lines are written, so no omission receipt of any kind accompanies this package. Citations: the retained excerpts carry no citation command at all, so no bibliography entry of the article is transcribed and no reference is redistributed. Reference adaptation: none was needed and none was made; every reference inside the delivered unit resolves inside it, so no unresolved reference or citation is produced. Printed slips kept literal: no printed word, symbol, hypothesis or number of the counted statement or of its proof was altered. Layout and class: the article's own class is \texttt{article} with packages including graphicx, amsmath, mathtools, amsthm, fancyhdr, amssymb, bbm, enumerate, float, mathrsfs, tcolorbox, appendix, listings, color; the wrapper is the lane's pinned \texttt{amsart} editorial wrapper at 10pt on A4, which restates the theorem-like environments the retained text needs and adds the article's own macro definitions that the retained text needs (none), with extra wrapper packages (bm, bbm, dsfont, xspace, cancel, mathtools, enumitem). The delivered numbering is the unit's own. Assets: no retained span contains a figure environment, a graphics-inclusion command, a PDF-page inclusion, a table or a TikZ picture; no image file is copied into this package, the package declares no asset dependency and no image file is redistributed with it. Licence: version arXiv:2604.12785v2 of this article is distributed under the Creative Commons Attribution 4.0 International licence, as recorded in the arXiv licence metadata for this exact version and shown on the abstract page https://arxiv.org/abs/2604.12785v2. A full rights notice is delivered at licenses/arxiv-2604.12785-CC-BY-4.0.txt. Characters: every retained excerpt is pure ASCII, so the delivered engine, XeLaTeX with its default Latin Modern text font, reports no missing character.
\begin{lemma}[Bochner Theorem]\label{Bochner lemma}
    Let $\phi:\mathbb{R}\rightarrow\mathbb{R}$ to be an even function with the condition that $\hat{\phi}(\xi)\geq0$. Then for any given  $x_1,x_2,\dots,x_n\in\mathbb{R}$, the following matrix
    \begin{equation}
        M=(\phi(x_i-x_j))_{n\times n}
    \end{equation}
    is symmetric and positive semi-definite.
\end{lemma}

\begin{proposition}\label{prop linear matrix}
    The matrix $A(\xi)$ is diagonalizable for all $\xi$ with $n$ eigenvalues $\lambda_1(\xi)\leq\lambda_2(\xi)\dots\leq\lambda_n(\xi)$ satisfying that:
    \begin{enumerate}[label=(\roman*), ref=(\roman*)]
        \item\label{item:positive-eigenvalues} For all $\xi\neq0$ we have \begin{equation}
            \lambda_k(\xi)>0,\quad \forall 1\leq k\leq n
        \end{equation} 
        and
        \begin{equation}
            \lambda_1(0)=\dots=\lambda_{n-1}(0)=0,\quad \lambda_n(0)=\frac{\rho_0-\rho_n}{2}.
        \end{equation}

        \item\label{item:asymptotic-eigenvalues}Asymptotic behavior: There exists $\epsilon,\alpha>0$ and a permutation \(\tau\) of \(\{1,\dots,n\}\) such that for all $0<\xi<\epsilon$ we have
        \begin{equation}
            \lambda_1(\xi)\geq\alpha\xi
        \end{equation}
        and
        \begin{equation}
            \lim_{\xi\rightarrow+\infty}\lambda_{\tau(k)}(\xi)=\frac{\rho_{k-1}-\rho_k}{2},\quad 1\leq k\leq n.
        \end{equation}
    \end{enumerate}
\end{proposition}

\begin{proof}
    We notice
    \begin{equation}
        A(\xi)=A_0(\xi)Q
    \end{equation}
    with
    \begin{equation}\label{def Q and A0}        A_0(\xi)=(e^{-|d_j-d_k||\xi|})_{n\times n},\quad Q=diag(\frac{\rho_{k-1}-\rho_k}{2})_{1\leq k\leq n } .  \end{equation}
    Then matrix $A$ is similar to $Q^{\frac{1}{2}}A_0Q^{\frac{1}{2}}$. Since $A_0$ is symmetric and $Q^{\frac{1}{2}}$ is a constant diagonal matrix with all components large than zero, we only need to prove \ref{item:positive-eigenvalues} and \ref{item:asymptotic-eigenvalues} for $A_0(\xi)$, with the asymptotic behavior replaced by
    \begin{equation}
        \lambda_n(0)=1,\quad\lim_{\xi\rightarrow+\infty}\lambda_{\tau(k)}(\xi)=1,\quad 1\leq k\leq n.
    \end{equation}
    We compute
    \begin{equation}
        \varphi(x)=e^{-|x|},\quad\widehat{\varphi}(\xi)= \frac{2}{1+\xi^2}.
    \end{equation}
    Then $A_0(\xi)$ is positive semi-definite  due to Lemma \ref{Bochner lemma} by noticing
    \begin{equation}
        A_0(\xi)=(\varphi(d_k\xi-d_j\xi))_{n\times n}.
    \end{equation}
    Furthermore when $\xi=0$, $A_0(\xi)$ is the all-ones matrix. Hence it has $n-1$ zero eigenvalues and one nonnegative eigenvalue equals $n$. When $\xi\neq 0$, it is straightforward to use upper triangularization to compute the determinant of $A_0(\xi)$ by replacing the $k$-th row $R_k$ with $R_k-e^{-(d_k-d_{k-1})|\xi|}R_{k-1}$ for $k=n,n-1,\dots,2$ successively to obtain
    \begin{equation}
        \text{det}(A_0(\xi))=\prod_{j=1}^{n-1}\left(1-e^{-2(d_{j+1}-d_j)|\xi|}\right)>0.
    \end{equation}
    This implies $A_0(\xi)$ is positive definite for $\xi\neq0$, therefore yields the first part of the proposition.

    When $\xi\rightarrow0^+$ we use the following expansion
    \begin{equation}
        A_0(\xi)=J+B\xi+O(\xi^2).
    \end{equation}
    Here
    \begin{equation}
        J=A_0(0)=\textbf{1}\textbf{1}^T,\quad B=(-|d_k-d_j|)_{n\times n},\quad \textbf{1}=(1,\dots,1)^T.
    \end{equation}
    We decompose the vector field as
    \begin{equation}
        \mathbb{R}^n=V\oplus V^{\perp}
    \end{equation}
    where $V$ denotes the eigenspace corresponding to the only nonzero eigenvalue of $J$
    \begin{equation}
        V=\text{span}\{\textbf{1}\}.
    \end{equation}
    Thus it is straightforward to decompose a unit vector $v=(v_1,v_2,\dots,v_n)^T\in V,\|v\|=1$ into
    \begin{equation}
        v=av'+bv'',\quad v'\in V,v''\in V^{\perp}
    \end{equation}
    with the normalization
    \begin{equation}
        a^2+b^2=1,\quad \|v'\|=1,\quad \|v''\|=1.
    \end{equation}
Direct computation shows
\begin{equation}
    v^TA_0v=a^2\left((v')^TJv'+O(\xi)\right)+\xi\left(2ab(v')^TBv''+b^2(v'')^TBv''\right)+O(\xi^2)
\end{equation}
Choosing a large  constant $C$ such that when $a\geq C|\xi|^{\frac{1}{2}}$,
\begin{equation}
    v^TA_0v\geq \frac{n}{2}C^2\xi-C\xi+O(\xi^2)\geq C\xi
\end{equation}
holds for all small $\xi$. Otherwise, if  $a\leq C|\xi|^{\frac{1}{2}}$, we observe   \begin{align}
        (v'')^TBv''&=\sum_{j,k=1}^n-|d_j-d_k|v''_jv''_k\\
        &=2\sum_{j=1}^{n-1}(d_{j+1}-d_j)\left(\sum_{k=1}^jv''_k\right)^2\geq 0
    \end{align}
    where we use $\sum_{j=1}^nv''_j=0$. Note the equality holds if and only if $v''=\textbf{0}$, then there exists $\mu>0$ which depends on $d_j,\rho_j$, such that
    \begin{equation}
        (v'')^TBv''\geq\mu. 
    \end{equation}
    This implies
    \begin{equation}
        v^TA_0v\geq a^2(n+O(\xi))+O(\xi^{\frac{3}{2}})+\frac{\mu}{2}\xi+O(\xi^2)\geq \frac{\mu}{4}\xi
    \end{equation}
    for sufficiently small $\xi$. The above argument shows by choosing $\alpha=\min\{C,\mu/4\}$ and sufficiently small  $\epsilon$ we have for all $0<\xi<\epsilon$,
    \begin{equation}
        \lambda_1(\xi)\geq\alpha\xi.
    \end{equation}

    For the asymptotic behavior of $A_0(\xi)$ at infinity we simply point out that $A_0(\xi)$ tends to identity matrix as $\xi\rightarrow\infty$. This completes the proof of proposition.
\end{proof}

\end{document}
